\chapter{Discussion}
This section interprets the numerical results in relation to the main goal of the thesis: to test whether numerical optimization can identify stable parameter regimes in interacting quantum-dot-superconductor chains that support Majorana-like low-energy physics and useful braiding behavior. It also discusses the physical implications of extending the braiding analysis beyond the ground-state manifold.

\section{Parameter Optimization Scheme}
For a simple two-dot system, the optimizer managed to approach the analytically known sweet spot, which provides a reassuring sanity check. In the non-interacting case, the gradient-based search found a parameter regime that matches the known sweet spot up to a numerical tolerance of approximately $10^{-3}$. This is a strong indication that the optimization procedure is working correctly, at least in the non-interacting limit.

For weakly interacting two-dot systems, the optimizer found parameter sets that are also close to the expected interacting sweet-spot conditions, where the ground-state degeneracy should be preserved. In particular, the optimized onsite energies remain close to $\epsilon_1,\epsilon_2 \approx -U/2$, with $\Delta \sim t$. This is consistent with the known interacting sweet-spot condition for two-dot systems, and suggests that the optimization is correctly identifying parameter regimes that support Majorana-like physics in the presence of weak interactions. For stronger interactions, the optimizer had more difficulty finding low-loss configurations. It also tended to push $\Delta > t$, while the relation between $\epsilon_i$ and $U$ remained close to the expected sweet-spot condition. This suggests that the optimizer may be reaching the physical limitations of the constrained parameter space, since $t$ was fixed to 1 throughout the optimization. When the interaction strength becomes too large compared with the hopping scale, the system may no longer be able to maintain even-odd degeneracy across its relevant energy levels.

Beyond the two-dot sweet spot, the analytical solution rapidly becomes more complicated. Without assuming symmetry or homogeneity, the number of parameters grows with the system size, and solving the sweet-spot conditions analytically becomes increasingly difficult. This makes numerical optimization a well-motivated approach for larger systems. For the three- and four-dot systems, the optimizer found parameter regimes that are not close to a simple homogeneous sweet-spot picture. Instead, the optimized parameters show a clear trend toward strong inhomogeneity, with large-magnitude onsite potentials in the middle of the chain and, for the four-dot system, enhanced pairing amplitudes near the center. This suggests that the optimizer is finding nontrivial parameter regimes that support Majorana-like physics, but that these regimes do not correspond to simple homogeneous configurations.

The parameter searches across all system sizes show that strong interactions make it substantially harder to find low-loss configurations. This is consistent with the expectation that strong interactions can destabilize Majorana-like physics, and it suggests that the optimizer is reflecting a real physical constraint rather than merely failing numerically.
However, at lower interaction strengths, especially in the three-dot system, the optimizer was able to find low-loss and seemingly stable configurations. This indicates that there are still weakly interacting parameter regimes that can support Majorana-like physics, and that the optimizer is capable of identifying these regimes.
For the four-dot system, the optimizer struggled more to find low-loss configurations even at weak interactions. This could reflect limitations of the optimization strategy, but it could also indicate that the parameter space becomes more complex and potentially more constrained as the system size increases. Since the optimization budget was kept fixed across all system sizes, the larger systems may not have been explored as thoroughly as the smaller ones. It would therefore be interesting to test whether increasing the number of iterations or using a more global search strategy could help the optimizer find better configurations for the four-dot system, especially in the weak-interaction regime.

One of the more interesting findings for chains longer than two dots was the emergence of large-magnitude onsite potentials in the middle of the chain. This was not expected from the simple sweet-spot picture. It suggests that the optimizer is finding nontrivial inhomogeneous configurations that may help support Majorana-like physics in interacting systems. This could have implications for the design of experimental devices, since it suggests that controlled inhomogeneity may play a stabilizing role rather than simply acting as an imperfection.

\section{Braiding Analysis}
The braiding results are mainly split into two parts. In the first construction, the Majorana operators are built from energy-level pairs of the optimized system. In the second construction, the braiding operators are built from projected local microscopic operators.

The main difference between the two constructions is what they test. The energy-level construction tests whether the spectrum supports a clean effective exchange channel that resembles the ideal Majorana model. A poor result in this construction would indicate that the relevant energy levels are not well isolated, or that the retained manifold contains states that do not behave like the intended Majorana degrees of freedom. The local-operator construction is more microscopic: it tests whether the same exchange channel can be accessed by projected local operators. A larger local-operator error can therefore indicate imperfect localization, hybridization with other states, target-basis mismatch, or sensitivity to the numerical tolerance of the optimized parameters.

The first part of the analysis tests whether braiding can be performed in the ground-state manifold using the Majorana modes generated by the optimized parameters. The results show that the ground-state manifold is well isolated, and that the projected junction operators closely match the ideal Majorana bilinears. This is a strong indication that the optimized parameters support Majorana-like physics in the ground-state manifold, and that the braiding protocol can be successfully implemented in this regime.

When higher-energy levels are included in the projection, the results show a clear separation between the weakly and strongly interacting cases. In the non-interacting case, the state-based Majorana operators remain close to the ideal Majorana operators across all projection spaces, and the projected junction braid remains close to the ideal control. With weak interactions included, the state-based Majorana construction shows a larger deviation, but the error peaks at about $10^{-4}$, which is still very good agreement. For strong interactions, however, the braid clearly begins to break down when more than the ground-state manifold is included. Interestingly, all three interaction strengths show a successful braid in the ground-state manifold, but the strong-interaction case shows a rapid increase in error once the projection space is enlarged. This suggests that the optimized parameters can support Majorana-like physics in the ground-state manifold even in the presence of strong interactions, while the higher-energy spectrum becomes more complicated and less well isolated.

The main visible difference between the local-operator braid and the energy-level braid is that the local-operator braid shows a larger error even in the non-interacting case, approaching the weak-interaction errors. This could indicate that the local-operator construction is more sensitive to the spatial structure of the system, and that even in the non-interacting case there may be small deviations from perfectly localized Majorana modes that are not visible in the energy-level construction. It could also be a consequence of numerical tolerance in the optimized parameters, which may affect the projected local operators more strongly than the state-based construction. This again highlights the importance of interpreting the results as evidence for approximate Majorana-like behavior, rather than as proof of exact topological protection.


\section{Limitations and Outlook}

While the gradient method provided insight into the structure and stability of the parameter regimes for the n-dot PMM systems, it is important to note a downside of the method. Since the gradient method provides a parameter regime optimized to minimize the loss function, it does not provide a perfect analytical sweetspot. Even the best parameter regimes found, are off by some numerical tolerance, which will transfer across to later diagnostics, as well as the braiding results as they are based on the parameters in use. Since the optimizer finds a numerical minimum which at best, minimizes down to a numerical error of around $10^{-3}$, a tolerance of this order should be expected in the later diagnostics, and in the braiding results. This means that the results should be interpreted as evidence for approximate and practically useful Majorana-like behavior, rather than as proof of exact topological protection.


The main limitation of the present study is that the optimization procedure does not locate exact sweet spots; it finds numerically favorable minima of the chosen loss function. In a complicated loss landscape, a value such as $10^{-3}$ may look essentially optimal even when a better solution exists elsewhere. Because all later diagnostics build on the optimized parameter sets, the braiding results should therefore be interpreted as evidence for approximate and practically useful Majorana-like behavior, not as proof of exact topological protection.

A second limitation is that the present braiding analysis is still highly controlled. The clean agreement found in the lowest manifold relies on projection to a small low-energy space and on an idealized, disorder-free setting. The study does not yet include disorder, quasiparticle poisoning, nonadiabatic driving errors, or explicit coupling to thermal environments. Those effects are precisely the kinds of ingredients that would determine whether the optimized configurations remain useful in a more realistic device-level setting.

Several natural extensions follow from this. On the optimization side, it would be valuable to test more global search strategies, alternative loss functions, and less restrictive parameter ans\"atze. On the braiding side, it would be interesting to study how disorder and noise affect the projected junction operators, and whether the weakly interacting configurations remain robust when the protocol is pushed beyond the ideal adiabatic limit. The higher-energy results also point to a more specific open question: to what extent can one deliberately optimize not only the ground-state manifold, but also nearby excited manifolds, in order to make the exchange protocol more tolerant to thermal occupation?
